% SGA 3, Editorial Notice and Introduction (Loop 1 English reconstruction)
% Source component: AaIntro-18fev11.pdf, local pages 1--6.
% Combined Polo--Gille reader: physical PDF pages 5--10.
% OCR witness: ocr_fo/AaIntro/full/AaIntro-18fev11/AaIntro-18fev11.md.
% The PDF controls structure, mathematical notation, and footnotes.

\begingroup
\newcommand{\SGAIIIIntroFootnote}[2]{%
  \begingroup
  \renewcommand{\thefootnote}{#1}%
  \footnote{#2}%
  \endgroup
}

\chapter*{编者告示}
\addcontentsline{toc}{chapter}{编者告示}
\markboth{编者告示}{编者告示}

这里呈现的是原研讨班讲义的略作修订的版本, 其宗旨与内容见导言.
修订工作主要包括改正排印错误, 在脚注中增添少量补充说明或参考文献,
将本书按现有形式分为三卷, 每卷附详细目录和记号索引, 并在第~3~卷末增设术语索引.
此外, J.-E. Bertin 撰写的第~VI\textsubscript{B}~篇由作者作了部分重写,
尤其是\SGARefLink{sga3:intro:section:5}{第~5~节}和第~10~节,
因而某些对该篇的引用与原版本中的引用有所不同. 本卷卷首列有本研讨班各篇的篇目.

自本研讨班讲义第一版问世以来, \emph{\'El\'ements de G\'eom\'etrie Alg\'ebrique}
第~IV~章已全部出版, 因而讲义中的某些段落已无必要. 我们有时在脚注中指明
EGA~IV 的有关出处, 读者据此可以略过这些段落.

关于另一种系统使用概形语言的代数群论述, 可参阅 M. Demazure 和 P. Gabriel 的
\emph{Groupes Alg\'ebriques} (North-Holland \& Masson et Cie).
该书与本研讨班讲义不同, 不以代数几何知识为先修要求, 而是包含所需的全部概形论预备知识;
因此, 它可以作为学习本研讨班讲义的入门读物. (该书还讨论了一些本研讨班未涉及的主题,
例如交换仿射代数群的 Dieudonn\'e 式结构理论, 见\emph{该书}第~V~章.)

\begin{flushright}
Bures-sur-Yvette, 1970 年 3 月
\end{flushright}

\clearpage
\thispagestyle{empty}
\mbox{}
\clearpage

\chapter*{导言}
\addcontentsline{toc}{chapter}{导言}
\markboth{导言}{导言}

\section*{\SGAInlineRefTarget{sga3:intro:section:1}1. 本研讨班有两个目的.}

一方面, 我们希望为一般群概形 (group scheme) 理论提供便于使用的基础.
第~I--IV~篇将提供概形与范畴语言方面不可或缺的预备练习. 为了获得一种自然契合
几何直观的语言, 并避免那些长此以往令人难以忍受的迂回说法, 我们总将一个
预概形~$S$~上的预概形~$X$ 与它所代表的函子
\[
  (\mathbf{Sch})_{/S}^{\circ}\longrightarrow \mathbf{Set}
\]
等同\SGAIIIIntroFootnote{\ensuremath{\ast}}{看来, 这种观点最早是 P. Cartier
在八九年前研究形式群理论时考虑到的; 遗憾的是, 他没有着手将这一观点确切表述并作
系统整理, 而它本来值得如此对待.}.
因此, 许多定义必须表述为适用于任意函子的形式, 而不论这些函子是否可表.
此外, 我们所用的函子几乎全都是“层” (关于“忠实平坦拟紧拓扑”).
第~IV~篇仅附带讨论群, 它勾勒出“局部化”和层的语言; 在涉及函子可表性的问题中,
这种语言也十分方便. 尤其是, 该篇将为后面出现的取商问题提供最便于使用的框架.

第~V~篇给出若干关于商的存在性的一般结果. 第~VI\textsubscript{A}~篇又研究这些结果,
讨论域上 (更一般地, Artin 环上) 的代数群对子群取商的问题
\SGAIIIIntroFootnote{\ensuremath{\ast\ast}}{关于取商, 尤其是对约化群取商的更深入研究,
可参阅 D. Mumford 的重要著作 \emph{Geometric Invariant Theory},
Ergebnisse der Mathematik, 第~34~卷, Springer, 1965.
我们指出, 该书的术语有一处重要之处与本书不同: 在特征为 $p>0$ 的域上,
Mumford 称为“约化”的群\textsuperscript{\SGARefLink{sga3:intro:section:1}{(1)}},
是本研讨班意义下的平滑乘法型群 (参见第~IX~篇). 可以合理地认为,
Mumford 采用“约化”的这种含义只是暂时的权宜之计, 因为当底不再是域时,
这一含义便失去了意义 (Mumford 早在其前言第二段中就出于其他理由作了大致这样的解释).
\textsuperscript{\SGARefLink{sga3:intro:section:1}{(1)}}\emph{编者注:}
见本导言末增补的说明.}.
该篇以及随后的第~VI\textsubscript{B}~篇还包含域上代数群所特有的各种初等结果,
这些结果在后文中经常用到.

第~VII~篇研究与底域的特征有关的某些现象, 特别是以适当的一般性建立
高度为~1~的根性群概形 (radicial group scheme) 与限制 $p$-Lie 代数之间的对应.

最后, 第~XVIII~篇给出了 Weil 关于代数群“双有理”定义的定理的概形论推广.

另一方面, 我们希望将仿射代数群的 Borel--Chevalley 结构理论推广到任意底预概形上的群.
事实上, 在撰写研讨班讲义的过程中, 我们发现这一理论的许多结果不需要仿射性假设.
当然, 最完备的结果是在“半单群概形”, 或更一般的“约化群概形”情形下获得的;
从第~XIX~篇起, 我们将只讨论这些群. Chevalley 本人已经在整数环上构造了
“T\^ohoku 群”, 本研讨班将重新研究这一构造. 主要的唯一性定理
\SGAIIIIntroFootnote{\SGARefLink{sga3:intro:section:2}{(2)}}{\emph{编者注:}
这里指推论~XXIII.5.6, 它容易由分裂约化群的唯一性定理推出
(参见 XXIII, \SGARefLink{sga3:XXIII:4.1}{定理~4.1}和\SGARefLink{sga3:XXIII:5.2}{推论~5.2}),
因为每个半单 $S$-群都是某个 T\^ohoku 群的“扭曲形式” (关于平展 (\(\acute{e}\)tale) 拓扑)
(参见 XXII, \SGARefLink{sga3:xxii:cor:2.3}{推论~2.3}).}
简洁地刻画了这些 T\^ohoku 群在底预概形~$S$~上的“扭曲”变体:
它们恰好是~$S$~上的仿射平滑群概形, 其几何纤维是通常意义下的连通半单群
\SGAIIIIntroFootnote{\ensuremath{\ast\ast\ast}}{这是 M. Demazure 博士论文
\emph{Sch\'emas en groupes r\'eductifs}, Bull. Soc. Math. France \textbf{93}
(1965), 369--413 的主要结果.}%
\textsuperscript{\SGARefLink{sga3:intro:section:2}{(2)}}.

\section*{\SGAInlineRefTarget{sga3:intro:section:2}2.}

与熟悉的代数闭域上的理论一样, 所研究的群概形的子环面发挥着关键作用.
因此, 对环面 (torus), 更一般地对“乘法型群概形”的预备研究,
无论研究其本身, 还是将它作为给定群的乘法型子群来研究, 都占据本研讨班相当大的篇幅
(第~VIII--XII~篇). 它们出色的刚性 (在某些方面甚至强于 Abel 概形或半单群概形的刚性)
使其成为研究某些更一般的群的有力工具.

\section*{\SGAInlineRefTarget{sga3:intro:section:3}3.}

从第~XII~篇起 (除去前面已提到的第~XVIII~篇), 我们将经常使用代数闭域上的
仿射代数群理论. 读者可在 1956 年 Chevalley 研讨班讲义中找到这些内容,
尤其是其中的第~IV--IX~篇. 我们也会用到 Chevalley 研讨班讲义中后面讨论半单代数群
结构的各篇, 不过用得较少. 事实上, 我们将直接在概形框架中重新建立 Chevalley 理论;
读者将看到, 即使在底为域的情形下, 这种做法也使论述更加简单而精确.

\section*{\SGAInlineRefTarget{sga3:intro:section:4}4.}

当然, 本研讨班的主要目的是发展适用于任意底上的群概形研究的技术,
也就是主要用于研究代数群族的技术. 因而, 这种族的无穷小性质, 特别是底概形为
Artin 概形的情形, 发挥着重要作用. 即使研究域~$k$~上的代数群,
当该域不完全时, 尤其在非 Galois 情形下应用下降时, 也会出现这些性质.
在这一背景下得到的新结果中, 可以提到任意平滑代数群中极大环面和 Cartan 子群的存在性,
极大环面簇的有理性, 以及各种相关结果 (第~XIV~篇
\SGAIIIIntroFootnote{\SGARefLink{sga3:intro:section:3}{(3)}}{\emph{编者注:}
特别见\SGARefLink{sga3:xiv:1.1}{定理~1.1}和~\SGARefLink{sga3:xiv:6.1}{6.1}.});
另一个例子是半单群的“形式”与适当的半单代数群 (通常不连通) 下的齐性主丛之间的对应
(\SGARefLink{sga3:xxiv:1.16}{第~XXIV~篇, 1.16}--\SGARefLink{sga3:xxiv:1.20}{1.20}).
一般说来, 在不完全底域上工作所需的方法, 本质上就是在任意底预概形上所用的方法,
因而超出了经典代数几何的框架.

\section*{\SGAInlineRefTarget{sga3:intro:section:5}5.}

我们认为没有必要在书面讲义各篇的开头注明研讨班中相应口头报告的日期.
这里只说明, 油印各篇 (I--XXVI) 的次序确实与口头报告的次序一致.
此外, 最终书面文本有时是在口头报告之后相当久才写成的, 而且往往与报告有很大差别:
书面文本通常更详细、更完整 (如第~IV~篇和第~VII\textsubscript{B}~篇),
甚至更为一般 (如第~XII~篇和第~VII\textsubscript{B}~篇).
另一些书面讲义则没有相应的口头报告
(VI\textsubscript{B}, VII\textsubscript{A}, XV, XVI, XVII, 以及 XXVI 的大部分);
撰写这些讲义并将其插入油印研讨班讲义中, 或是为了给其他各篇提供便于使用的参考资料
(尤其是 VI\textsubscript{B} 的情形), 或是因为它们是已经展开的概念和技术的自然延续.
因此, 为了学习研讨班的其余部分, 特别是约化群概形部分,
不必阅读第~VII\textsubscript{A}, VII\textsubscript{B}, XV, XVI 和 XVII~篇.

\section*{\SGAInlineRefTarget{sga3:intro:section:6}6.}

在概形论方面, 我们主要使用 EGA~I 中阐述的概形的一般语言;
SGA~1, I--V 中阐述的平坦、平展 (\(\acute{e}\)tale) 及平滑态射的概念;
以及 SGA~1, VIII 中的忠实平坦下降理论. 我们尽可能避免施加不必要的 Noether 假设.
相应地, 这就要求我们把通常的“有限型”假设替换为“有限呈示”假设.
关于有限呈示态射的概念, 读者应参阅 EGA~IV, 1.4 和~1.6.
SGA~1, I--IV 中的结果大多在 Noether 情形下陈述, EGA~IV 将在一般情形下建立这些结果
\SGAIIIIntroFootnote{\ensuremath{\ast\ast\ast\ast}}{自本导言写成以来,
EGA~IV 的全部四个部分 (第~1--21~节) 均已出版.};
该章还将详细介绍把涉及有限呈示假设的某些类型的命题化归到 Noether 情形的标准方法
(EGA~IV, 第~8, 9 和~11~节). 不愿预先承认 EGA~IV 中这些结果的读者,
可以假定底预概形局部 Noether, 从而简化某些陈述或其证明.
不过, 当证明要从~$\widehat A$~下降到~$A$~时, 这种做法会带来困难,
其中 $\widehat A$ 是 Noether 局部环~$A$~的完备化, 因为这种方法引入的环
\[
  \widehat A\otimes_A\widehat A.
\]
一般不是 Noether 环.

\section*{\SGAInlineRefTarget{sga3:intro:section:7}7.}

引用遵循通常的十进制编号体系: 引用 5.7.11 表示同一篇中具有该编号的命题
(或引理、定义等); 在 XVII~7.8 这样的引用中, 罗马数字表示篇号.
对于经常引用的文献, 我们采用下列缩写:

\begin{description}
\item[\emph{Bible}] Chevalley 研讨班, \emph{Groupes de Lie alg\'ebriques},
  1956/58.
\item[EGA $X$, $x.y.z$] J. Dieudonn\'e 和 A. Grothendieck,
  \emph{\'El\'ements de G\'eom\'etrie Alg\'ebrique}, 第~$X$~章,
  编号为 $x.y.z$ 的陈述 (或第 $x.y$ 小节, 等等).
\item[SGA $n$, $X$ $y.z$] \emph{S\'eminaire de G\'eom\'etrie Alg\'ebrique du
  Bois-Marie}, 第~$n$~年, 第~$X$~篇中编号为 $y.z$ 的陈述.
\item[TDTE] A. Grothendieck, \emph{Techniques de descente et Th\'eor\`emes
  d'existence en G\'eom\'etrie Alg\'ebrique}, 1959 至 1962 年间在 Bourbaki
  研讨班所作的各篇报告.
\end{description}

\begin{quote}\small
\noindent\textsuperscript{\SGARefLink{sga3:intro:section:1}{(1)}}\emph{编者注.}
关于对约化群取商的理论, 自 A. Grothendieck 写下脚注~$(\ast\ast)$~以来,
情况已经有了很大变化. 对任意域~$k$~上的仿射代数群, 恰当的概念是 Mumford
引入的“几何约化”群. (另一方面, 若~$G$~的每个有理表示都完全可约,
则称 $G$ 为“线性约化”的; 但是, 如脚注~$(\ast\ast)$~所指出的,
当 $\operatorname{char}(k)>0$ 时, 这一条件过于苛刻.)
由 M. Nagata 和 W. J. Haboush 的结果 ([Na64], [NM64], [Ha75]),
几何约化 $k$-群恰好就是本研讨班意义下的约化 $k$-群.
关于这些内容, 可参阅 Mumford 著作的后续版本 ([MF82]).
C. S. Seshadri 将几何约化性的概念推广到任意底的情形, 并研究了它对取商的推论
([Se77]); J.-L. Colliot-Th\'el\`ene 和 J.-J. Sansuc 的文章 [CTS79]
中载有 M. Raynaud 的补充结果. 最后, 关于对代数群取商, 或更一般地对群胚取商
(参见本研讨班第~V~篇) 的较新进展, 我们推荐 J. Koll\'ar 的文章 [Ko97]
以及 S. Keel 和 S. Mori 的文章 [KM97].
\end{quote}

\section*{参考文献}

\begin{description}
\item[{[CTS79]}] J.-L. Colliot-Th\'el\`ene 和 J.-J. Sansuc,
  \emph{二次纤维丛与实连通分支〔Fibr\'es quadratiques et composantes connexes r\'eelles〕}, Math. Ann.
  \textbf{244} (1979), 105--134.
\item[{[Ha75]}] W. J. Haboush, \emph{约化群是几何约化的〔Reductive groups are geometrically
  reductive〕}, Ann. of Math. \textbf{102} (1975), 第~1~期, 67--83.
\item[{[KM97]}] S. Keel 和 S. Mori, \emph{群胚的商〔Quotient by groupoids〕}, Ann. of Math.
  \textbf{145} (1997), 第~1~期, 193--213.
\item[{[Ko97]}] J. Koll\'ar, \emph{模代数群的商空间〔Quotient spaces modulo algebraic groups〕}, Ann.
  of Math. \textbf{145} (1997), 第~1~期, 33--79.
\item[{[MF82]}] D. Mumford 和 J. Fogarty, \emph{几何不变量理论〔Geometric invariant theory〕},
  第~2~版, Springer-Verlag, 1982; 第~3~版, 与 F. Kirwan 合著, 1994.
\item[{[Na64]}] M. Nagata, \emph{仿射环中群的不变量〔Invariants of a group in an affine ring〕},
  J. Math. Kyoto Univ. \textbf{3} (1964), 第~3~期, 369--377.
\item[{[NM64]}] M. Nagata 和 T. Miyata, \emph{关于半约化群的注记〔Note on semi-reductive groups〕},
  J. Math. Kyoto Univ. \textbf{3} (1964), 第~3~期, 379--382.
\item[{[Se77]}] C. S. Seshadri, \emph{任意底上的几何约化性〔Geometric reductivity over an arbitrary
  base〕}, Adv. Math. \textbf{26} (1977), 第~3~期, 225--274.
\end{description}

\endgroup


% 书目补译核对：8个著作／论文／论文集题名已附中文，原题完整保留于〔〕；人名、刊名、出版机构、卷期页码与引用键保留。
